%%%%%%%%%%  scan idx 139  |  folio 126  %%%%%%%%%%
\SGAPrintPage{126}%
\sgaStableTarget{sga6.e01.rem.4-12-2}\textbf{Remark 4.12.2.} Let \(\underline{C}\) be an abelian category (which can be regarded as fibered over a punctual site), and let \(\underline{C}_0\) be a strictly full additive subcategory. Suppose that \(\underline{C}_0\) is stable under kernels of epimorphisms and \uline{quasi-liftable} in \(\underline{C}\). Then an acyclic object of \(\Kpx^b(\underline{C}_0)\) is in general not zero, so there is no hope that the functor
\(\Kpx^b(\underline{C}_0)\longrightarrow \Dpx(\underline{C})\)
is faithful. However, when \(\Kpx^{b,\varnothing}(\underline{C}_0)\) denotes the thick subcategory of \(\Kpx^b(\underline{C}_0)\) formed by the acyclic objects (in \(\underline{C}\)), the canonical functor
\[
  \Kpx^b(\underline{C}_0)/\Kpx^{b,\varnothing}(\underline{C}_0)
  \longrightarrow \Dpx(\underline{C})
\]
is fully faithful.

Hints: We already know (2.7) that the functor
\[
  \Kpx^-(\underline{C}_0)/\Kpx^{-,\varnothing}(\underline{C}_0)
  \longrightarrow \Dpx^-(\underline{C})
\]
is fully faithful. Show, for example using the criterion ([V] 4.2.b)(ii)'), that the functor
\[
  \Kpx^b(\underline{C}_0)/\Kpx^{b,\varnothing}(\underline{C}_0)
  \longrightarrow
  \Kpx^-(\underline{C}_0)/\Kpx^{-,\varnothing}(\underline{C}_0)
\]
is fully faithful.

\sgaStableTarget{sga6.e01.lem.4-13}\textbf{Lemma 4.13.} \uline{Let} \(F\in\Ob\Dpx(\underline{C}_S)\)\uline{,} \(a,b,c\in\mathbb Z\) \uline{be such that} \(a\leq b\) \uline{and} \(a\leq c\)\uline{. If} \(F\) \uline{has perfect amplitude contained in} \([a,c]\) (4.7) \uline{and} \(\Ho^i(F)=0\) \uline{for} \(i>b\)\uline{, then} \(F\) \uline{has perfect amplitude contained in} \([a,d]\)\uline{, where} \(d=\inf(b,c)\)\uline{.}

\uline{Proof.} We can suppose that \(b\leq c\), and that \(F\) is strictly perfect with \(F^i=0\) for \(i\notin[a,c]\). Then \(F\) is isomorphic to its truncation \(F'\) defined by
%%%%%%%%%%  scan idx 140  |  folio 127  %%%%%%%%%%
\SGAPrintPage{127}%
\(F'^i=F^i\) for \(i<b\), \(F'^b=Z^b(F)\), \(F'^i=0\) for \(i>b\). But \(Z^b(F)\) admits a finite right resolution by objects of \(\underline{C}_0\), so it is locally an object of \(\underline{C}_0\), whence the assertion.

\sgaStableTarget{sga6.e01.cor.4-14}\textbf{Corollary 4.14.} \uline{Let} \(F\in\Ob\underline{C}_S\) \uline{and} \(n\in\mathbb N\)\uline{. Equivalent conditions}:
\begin{enumerate}[label=(\roman*)]
\item \(\underline{\perfamp}(F)\subset[-n,0]\);
\item \(F\) \uline{locally admits a left resolution of length} \(n\) \uline{by objects of} \(\underline{C}_0\)\uline{, i.e.\ the set of} \(X\) \uline{over} \(S\) \uline{for which there exists an exact sequence}
\[
  0\longrightarrow L^{-n}\longrightarrow L^{-n+1}\longrightarrow \ldots
  \longrightarrow L^0\longrightarrow F_X\longrightarrow 0,
\]
\uline{with} \(L^i\in\Ob\underline{C}_{0X}\) \uline{for every} \(i\), \uline{is a refinement of} \(S\).
\end{enumerate}

\uline{Proof.} This is an immediate consequence of (4.7) and (4.13).

\sgaStableTarget{sga6.e01.def.4-14-1}\textbf{Definition 4.14.1.} \uline{We shall say that} \(F\) \uline{has perfect dimension} \(\leq n\)\uline{, and we shall write} \(\perfdim(F)\leq n\)\uline{, when} \(F\) \uline{satisfies the equivalent conditions of} (4.14)\uline{. We shall call the smallest integer} \(n\) \uline{such that} \(\underline{\perfdim}(F)\leq n\) \uline{the} perfect dimension \uline{of} \(F\) \uline{(when there is no such integer, we shall say that} \(F\) \uline{has infinite perfect dimension).}

When \(\underline{S}\) is the punctual site and \(\underline{C}_0\) is the full subcategory of \(\underline{C}\) formed by the projective objects, perfect dimension is the usual projective dimension.

\sgaStableTarget{sga6.e01.prop.4-15}\textbf{Proposition 4.15.} \uline{Let} \(a,b,n\in\mathbb Z\) \uline{(}\(a\leq b\)\uline{,} \(n\geq0\)\uline{), and let} \(F\in\Ob\Dpx(\underline{C})\) \uline{be such that} \(F^i=0\) \uline{(resp.}\ \(\Ho^i(F)=0\)\uline{) for} \(i\notin[a,b]\)\uline{. If, for every} \(i\in[a,b]\)\uline{,} \(F^i\) \uline{(resp.}\ \(\Ho^i(F)\)\uline{) has perfect dimension} \(\leq n+(i-a)\)\uline{, then} \(F\) \uline{has perfect amplitude contained in} \([a-n,b]\)\uline{. Consequently, if for every} \(i\in[a,b]\)\uline{,} \(F^i\) \uline{(resp.}\ \(\Ho^i(F)\)\uline{) is perfect,} \(F\) \uline{is too.}

%%%%%%%%%%  scan idx 141  |  folio 128  %%%%%%%%%%
\SGAPrintPage{128}%
\uline{Proof.} Proceed by induction on \(b\) (with \(a\) and \(n\) fixed), truncating \(F\) and applying (4.11) (cf.\ the proof of (2.11)).

\sgaStableTarget{sga6.e01.lem.4-16}\textbf{Lemma 4.16.} \uline{Let} \(a,b\in\mathbb Z\) \uline{(}\(a\leq b\)\uline{), and} \(F\in\Ob\Dpx(\underline{C})\) \uline{be such that} \(F^i=0\) \uline{for} \(i\notin[a,b]\)\uline{. Assume that} \(\underline{\perfamp}(F)\subset[a,b]\)\uline{, and that} \(F^i\in\Ob\underline{C}_0\) \uline{for} \(i>a\)\uline{. Then} \(F^a\) \uline{is locally an object of} \(\underline{C}_0\)\uline{.}

\uline{Proof.} We have an exact sequence (obtained by truncating \(F\))
\[
  0\longrightarrow F'\longrightarrow F\longrightarrow F^a[-a]\longrightarrow 0,
\]
with \(\perfamp(F')\subset[a+1,b+1]\). From (4.11) and (4.13) we deduce that \(\perfdim(F^a)=0\), i.e.\ that \(F^a\) is locally an object of \(\underline{C}_0\).

\sgaStableTarget{sga6.e01.prop.4-17}\textbf{Proposition 4.17.} \uline{Let} \(F',F''\in\Ob\Dpx(\underline{C}_S)\)\uline{, and} \(a,b\in\mathbb Z\)\uline{, with} \(a\leq b\)\uline{. Consider the conditions}:
\begin{enumerate}[label=(\roman*)]
\item \(\perfamp(F')\subset[a,b]\) and \(\perfamp(F'')\subset[a,b]\);
\item \(\perfamp(F'+F'')\subset[a,b]\).
\end{enumerate}
\uline{We have} (i) \(\Rightarrow\) (ii)\uline{, and, when} \(\underline{C}_0\) \uline{is locally stable under direct factors (}i.e.\ \uline{the relation} \(E'+E''\in\Ob\underline{C}_0\)\uline{, for} \(E',E''\in\Ob\underline{C}_X\) \uline{and} \(X\in\Ob \underline{S}\)\uline{, implies that locally} \(E'\) \uline{and} \(E''\) \uline{are objects of} \(\underline{C}_0\)\uline{), we have} (i) \(\Leftrightarrow\) (ii).

\uline{Proof.} The implication (i) \(\Rightarrow\) (ii) is a special case of (4.11). Let us prove the converse. Put \(F=F'+F''\). Since \(F\) is pseudo-coherent, \(F'\) and \(F''\) are too (2.12). On the other hand, the relation \(\Ho^i(F)=0\) for \(i\notin[a,b]\) implies \(\Ho^i(F')=\Ho^i(F'')=0\) for \(i\notin[a,b]\). By (2.10 a)), we can therefore suppose (after localizing) that \(F'^i\) (resp.\ \(F''^i\)) \(\in\Ob\underline{C}_0\) for \(i\geq a+1\) and \(F'^i=F''^i=0\) for \(i>b\), then, by truncating, that
%%%%%%%%%%  scan idx 142  |  folio 129  %%%%%%%%%%
\SGAPrintPage{129}%
\(F'^i=F''^i=0\) for \(i<a\). Since \(\perfamp(F)\subset[a,b]\), we deduce from (4.16) that \(F^a\) is locally an object of \(\underline{C}_0\). But \(F^a=F'^a+F''^a\), so \(F'^a\) and \(F''^a\) are locally objects of \(\underline{C}_0\), since \(\underline{C}_0\) is locally stable under direct factors.

\sgaStableTarget{sga6.e01.prop.4-18}\textbf{Proposition 4.18.} \uline{Let} \((\underline{S}',\underline{C}',\underline{C}'_0)\) \uline{be a triple satisfying the same hypotheses} (4.0) \uline{as} \((\underline{S},\underline{C},\underline{C}_0)\)\uline{. Let} \(f:\underline{S}'\to \underline{S}\) \uline{be a morphism of sites (defined by} \(f^{-1}:\underline{S}\to \underline{S}'\)\uline{), and let}
\(u:\Dpx^*(\underline{C})\longrightarrow \Dpx^*(\underline{C}')\qquad (*=-,+,\varnothing)\)
\uline{be a functor over} \(f\) (cf.\ (2.13)). \uline{If} \(u\) \uline{takes objects of} \(\underline{C}_0\) \uline{to} \(\underline{C}'_0\)\uline{-perfect complexes, then} \(u\) \uline{defines a functor}
\(\Dpx^*(\underline{C})_{\underline{C}_0\text{-}\mathrm{parf}}\longrightarrow\Dpx^*(\underline{C}')_{\underline{C}'_0\text{-}\mathrm{parf}}\).

\uline{Proof.} Let \(F\in\Ob\Dpx^*(\underline{C})_{\underline{C}_0\text{-}\mathrm{parf}}\); let us show that \(u(F)\in\Ob\Dpx^*(\underline{C}')_{\underline{C}'_0\text{-}\mathrm{parf}}\). Since the question is local, we can suppose that \(F\) is strictly perfect. We need only use dévissage, taking into account that \(\Dpx^*(\underline{C}')_{\underline{C}'_0\text{-}\mathrm{parf}}\) is a triangulated subcategory of \(\Dpx^*(\underline{C}')\) (4.10).

\sgaStableTarget{sga6.e01.rem.4-18-1}\textbf{Remark 4.18.1.} Suppose that there exist \(m,n\in\mathbb Z\) (\(m\leq n\)) such that \(F\in\Ob\underline{C}_0\) implies \(\underline{C}'_0\text{-}\perfamp(u(F))\subset[m,n]\). Then, for \(F\in\Ob\Dpx^*(\underline{C})\), we have
\[
  \underline{C}_0\text{-}\perfamp(F)
  \subset[a,b]
  \quad\Rightarrow\quad
  \underline{C}'_0\text{-}\perfamp(u(F))
  \subset[a+m,b+n].
\]
The proof is left as an exercise to the reader (use (4.11)).

\sgaStableTarget{sga6.e01.cor.4-19}\textbf{Corollary 4.19.} \uline{Under the hypotheses of} (2.16), \uline{if} \(E'\in\Ob\Dpx^-({}_{A'} \underline{S}'_A)\) \uline{is perfect as an object of} \(\Dpx^-({}_{A'} \underline{S}')\) \uline{and} \(E\in\Ob\Dpx^-({}_A \underline{S})\) \uline{is perfect, then}
\(E'\overset{\mathrm L}{\otimes}_A E\) \uline{is perfect}.

\uline{Proof.} We can suppose that \(E'\in\Ob\Dpx^-({}_{A'} S'_A)\) (\(S'\)= final object of \(\underline{S}'\)), and apply (4.18) to
\(u=E'\overset{\mathrm L}{\otimes}_A:\Dpx^-({}_A \underline{S})\longrightarrow \Dpx^-({}_{A'} \underline{S}')\).
In particular:

%%%%%%%%%%  scan idx 143  |  folio 130  %%%%%%%%%%
\SGAPrintPage{130}%
\sgaStableTarget{sga6.e01.cor.4-19-1}\textbf{Corollary 4.19.1.} a) \uline{The functor}
\(\mathbb L f^*:\Dpx^-({}_A \underline{S})\longrightarrow \Dpx^-({}_{A'} \underline{S}')\)
\uline{induces a functor}
\[
  \mathbb L f^*:\Dpx^-({}_A \underline{S})_{\mathrm{parf}}
  \longrightarrow
  \Dpx^-({}_{A'} \underline{S}')_{\mathrm{parf}}.
\]

b) \uline{If} \(A\) \uline{is commutative, the functor}
\(\overset{\mathrm L}{\otimes}_A:\Dpx^-(\underline{S})\times_{\underline{S}}\Dpx^-(\underline{S})\longrightarrow \Dpx^-(\underline{S})\)
\uline{induces a functor}
\[
  \overset{\mathrm L}{\otimes}_A:
  \Dpx^-(\underline{S})_{\mathrm{parf}}\times_{\underline{S}}\Dpx^-(\underline{S})_{\mathrm{parf}}
  \longrightarrow \Dpx^-(\underline{S})_{\mathrm{parf}}.
\]

\sgaStableTarget{sga6.e01.rem.4-19-2}\textbf{Remark 4.19.2.} We can make (4.19) more precise in terms of perfect amplitude:
\[
  \perfamp(E)\subset[a,b]\quad\text{and}\quad
  \perfamp(E')\subset[a',b']
  \quad\Rightarrow\quad
  \perfamp(E'\overset{\mathrm L}{\otimes}_A E)\subset[a+a',b+b'].
\]
In particular:
\[
  \perfamp(E)\subset[a,b]
  \quad\Rightarrow\quad
  \perfamp(\mathbb L f^*(E))\subset[a,b].
\]

\sgaStableTarget{sga6.e01.sec.5}\section*{5. Finite Tor-dimension and perfection}

\sgaStableTarget{sga6.e01.para.5-0}5.0. The reader will surely have asked: what does a pseudo-coherent complex \(E\) lack in order to be perfect? Is it enough, for example, that \(E\) have bounded cohomology? The answer is no. Indeed, take for \(\underline{C}\) the category of modules over a commutative noetherian ring \(A\), and for \(\underline{C}_0\) the full subcategory of projective \(A\)-modules of finite type. Any \(A\)-module of finite type \(E\) is pseudo-coherent and has bounded cohomology, but \(E\) is perfect only when \(E\) also has finite projective dimension (or, equivalently, finite tor-dimension). This example suggests that the difference between a pseudo-coherent complex and a perfect complex lies in a question of cohomological dimension. This is what we shall now examine
